\documentclass[11pt]{article}

\usepackage{iftex}
\ifPDFTeX
  \usepackage[utf8]{inputenc}
  \usepackage[T1]{fontenc}
  \usepackage{lmodern}
\fi
\usepackage[british]{babel}
\usepackage[a4paper,margin=2.6cm]{geometry}
\usepackage{amsmath,amssymb,amsthm}
\usepackage{microtype}
\usepackage{booktabs}
\setlength{\emergencystretch}{1.5em}

\newtheorem{satz}{Theorem}[section]
\newtheorem{lemma}[satz]{Lemma}
\newtheorem{korollar}[satz]{Corollary}
\newtheorem{vermutung}[satz]{Conjecture}
\theoremstyle{definition}
\newtheorem{definition}[satz]{Definition}
\newtheorem{beispiel}[satz]{Example}
\theoremstyle{remark}
\newtheorem{bemerkung}[satz]{Remark}

\newcommand{\Hori}{\mathcal H}
\newcommand{\Struk}{\mathcal S}
\newcommand{\NN}{\mathbb N_0}
\newcommand{\val}{\operatorname{val}}
\newcommand{\ff}[2]{#1^{\underline{#2}}}
\newcommand{\KV}{K_V}
\newcommand{\KS}{K_S}
\newcommand{\hor}{\operatorname{hor}}

\title{Horimetrik:\\
Value and Structure in a Positional Number System\\
with Finite Factorial Horizons}
\author{Christian Felix Bürckert\thanks{Email:
\texttt{christian@buerckert.eu}. Code, data, and an extended
exposition (in German and English) are available at
\texttt{github.com/christianbuerckert/horimetrik} and archived
at \texttt{doi:10.5281/zenodo.21400755}. Licence: CC\,BY\,4.0.}}
\date{14 July 2026 --- preprint, v0.16}

\begin{document}
\maketitle

\begin{abstract}
We study a mixed-radix positional number system whose bases
\emph{decrease} towards the most significant position: a digit
string $a_1\cdots a_n$ with $0\le a_i\le i$ denotes the value
$\sum_i a_i\,(n+1)!/(i+1)!$. Such a system is sealed to the left;
every number lives in a finite \emph{horizon} of size $(n+1)!$, and
a leading zero changes the value. Every element thereby carries two
identities: its \emph{value} and its \emph{structure} (a polynomial
with non-negative coefficients in the basis of falling factorials).
We show: (1)~The value- resp.\ structure-preserving embeddings of
the finite horizons generate two different direct limit spaces,
$\NN$ and the positive structure semiring. (2)~The space carries two
mirror-image commutative semirings; in both, the same distinguished
one acts as a projector onto the respective canonical copy. (3)~A
\emph{side-fidelity theorem}: no commutative semiring arithmetic
with value- resp.\ structure-carried operations can mix the two
identities, independently of the choice of horizon rules; the
obstruction is precisely the digit bound of the system. In addition
we prove an asymptotic law for the commutator of the two
compactifications, an extremal principle with exact growth laws of
the structure horizon under addition, multiplication, finite
difference, and composition (the composition closure of the
positive semiring itself is part of the result), and document nine
apparently unknown counting sequences. The text makes no claim of
novelty beyond a documented search finding.
\end{abstract}

\noindent\textbf{MSC 2020:} 05A15, 11A63, 16Y60.\quad
\textbf{Keywords:} mixed-radix number systems, falling factorials,
semirings, inversion sequences, integer sequences.

\section{The idea}

Ordinary positional systems let the most significant position count
farthest. This text studies the reverse convention. For fixed $n$,
let a digit string $a_1a_2\cdots a_n$ with $0\le a_i\le i$ denote
the value
\[
V_n(a)=\sum_{i=1}^n a_i\,\frac{(n+1)!}{(i+1)!}\,;
\]
the leftmost position is binary, yet carries the largest weight.
Two consequences: the system is \emph{sealed to the left} --- below
base $2$ there is nothing, with $n$ digits exactly $(n+1)!$ values
are representable, and more room requires a new coordinate system.
And \emph{leading zeros are semantic}: $12\mapsto012\mapsto0012$
denotes $5,6,7$ in turn.

Every number here thus carries two things: a value and a shape. The
mathematics of this dual character is the subject of this text. The
idea arose from the question of what happens if one reverses the
exponent ordering of a number representation; the finiteness of the
horizons is not a design goal but a necessity.

\section{The Hori space}

\begin{definition}
For $n\ge0$ let $H_n=\{(a_1,\dots,a_n):0\le a_i\le i\}$, so
$|H_n|=(n+1)!$, with the value map $V_n$ as above and the encoding
$E_n=V_n^{-1}$. Further let
$\mu(x)=\min\{n:x<(n+1)!\}$ be the \emph{minimal horizon} of
$x\in\NN$.
\end{definition}

\begin{satz}\label{satz:bij}
$V_n\colon H_n\to\{0,\dots,(n+1)!-1\}$ is bijective.
\end{satz}

\begin{proof}
Induction on $n$ via
$V_n(a)=(n+1)V_{n-1}(a_1,\dots,a_{n-1})+a_n$ and $0\le a_n\le n$:
division with remainder by $n+1$ yields uniqueness; the range is
$[0,(n+1)!-1]$, and both sides have cardinality $(n+1)!$.
\end{proof}

\begin{definition}
With $\ff Xd=X(X-1)\cdots(X-d+1)$ let
$\Struk=\bigl\{\sum_d c_d\ff Xd:c_d\in\NN\bigr\}$ and
$\sigma\bigl(\sum c_d\ff Xd\bigr)=\max_{c_d>0}(d+c_d)$,
$\sigma(0)=0$. The \emph{structure polynomial} of $a\in H_n$ is
$P_a=\sum_{d=0}^{n-1}a_{n-d}\ff Xd$; we have $V_n(a)=P_a(n+1)$, and
$P\in\Struk$ is representable in $H_n$ if and only if
$\sigma(P)\le n$. The \emph{Hori space} is
$\Hori=\{(n,P):P\in\Struk,\ n\ge\sigma(P)\}$, equivalently
$\{(n,x):n\ge\mu(x)\}$.
\end{definition}

$\Struk$ is closed under addition and multiplication (the product
formula of the falling factorials has non-negative structure
constants) and is therefore a commutative semiring.

\section{First main result: two limit spaces}

There are two natural ways to give an element more room: the
\emph{structure-preserving extension} $Z(n,P)=(n+1,P)$ (leading
zero; the shape stays, the value becomes $P(n+2)$) and the
\emph{value-preserving lift} $L(n,P)=(n+1,E_{n+1}(P(n+1)))$ (the
value stays, the shape is re-encoded). In the decimal system the
two coincide; here they do not.

\begin{satz}\label{satz:limits}
As directed systems,
\[
\varinjlim(H_n,Z_n)\cong\Struk,
\qquad
\varinjlim(H_n,L_n)\cong\NN .
\]
\end{satz}

\begin{proof}
Under $Z$, exactly those digit strings are identified that differ
by leading zeros; each class carries exactly one structure
polynomial, and every $P\in\Struk$ occurs (from $H_{\sigma(P)}$
onwards). Under $L$ we have $L_{m,k}\circ L_{n,m}=L_{n,k}$ because
$V\circ E=\mathrm{id}$, and two elements are identified precisely
when their values agree; every value occurs.
\end{proof}

The same finite spaces thus generate, depending on the notion of
consistency between the horizons, two different infinite worlds:
the values $\NN$ and the structures $\Struk$. Everything that
follows concerns the relationship between these two worlds.

\begin{bemerkung}[third extension; Cantor coordinate]
There is a third natural extension: \emph{appending} a zero on the
right, $R(a_1,\ldots,a_n)=(a_1,\ldots,a_n,0)$. Because
$V_n(a)=(n+1)!\cdot F(a)$ with the factorial fraction
$F(a)=\sum_i a_i/(i+1)!\in[0,1)$, the map $R$ is
\emph{fraction-preserving} and scales the value exactly by $n+2$;
the associated limit space is the set of finite Cantor factorial
fractions. The obvious objection ``essentially known since
Cantor'' thus applies exactly to this third, fraction-preserving
world -- not to the shape-preserving one (whose limit space is the
structure semiring) and not to the theorems on the interplay proved
below. Fraction, shape, and value are pairwise distinct invariants.
\end{bemerkung}

\section{Second main result: dual projector semirings}

The two \emph{compactifications}\footnote{Not meant topologically:
idempotent retractions onto the respective minimal horizon.}
$\KV(n,x)=(\mu(x),x)$ (keep the value, smallest home) and
$\KS(n,P)=(\sigma(P),P)$ (keep the shape, smallest home) are
idempotent and do not commute; the smallest example is the digit
string $010$ in $H_3$ (value 4).

\begin{satz}[value side]\label{satz:wertseite}
With
\[
(n,x)\oplus(m,y)=\bigl(\max(n,m,\mu(x{+}y)),\,x{+}y\bigr),
\qquad
(n,x)\otimes(m,y)=\bigl(\mu(xy),\,xy\bigr)
\]
$\Hori$ is a commutative semiring without a one; $(0,0)$ is neutral
and absorbing; the $\KV$-fixed points form a subsemiring isomorphic
to $(\NN,+,\cdot)$ with one $(1,1)$, and for all $h$ we have
\[
h\otimes(1,1)=\KV(h).
\]
\end{satz}

\begin{proof}
The core is the monotonicity of $\mu$. Associativity of $\oplus$:
both bracketings yield
$\max(n,m,k,\mu(\text{partial sum}),\mu(x{+}y{+}z))$, and the inner
$\mu$-terms are absorbed by $\mu(x{+}y{+}z)$. Distributivity:
$\max(\mu(xy),\mu(xz),\mu(xy{+}xz))=\mu(x(y{+}z))$, again
monotonicity. Neutrality of $(0,0)$ follows from $n\ge\mu(x)$. An
$\otimes$-neutral element would have to force $\mu(x)=n$ for all
elements, which fails at $(1,0)$. Finally
$(n,x)\otimes(1,1)=(\mu(x),x)=\KV(n,x)$.
\end{proof}

\begin{lemma}\label{lem:sigmamono}
$\sigma(P+Q)\ge\max(\sigma(P),\sigma(Q))$.
\end{lemma}

\begin{proof}
At the maximal position $d^*$ of $P$ we have
$c_{d^*}(P{+}Q)\ge c_{d^*}(P)>0$.
\end{proof}

\begin{satz}[structure side, duality]\label{satz:dualitaet}
With
\[
(n,P)\boxplus(m,Q)=\bigl(\max(n,m,\sigma(P{+}Q)),\,P{+}Q\bigr),
\qquad
(n,P)\boxtimes(m,Q)=\bigl(\sigma(PQ),\,PQ\bigr)
\]
$\Hori$ is likewise a commutative semiring without a one; the
$\KS$-fixed points form a subsemiring isomorphic to $\Struk$;
$\KS$ is a homomorphism for both operations, and for all $h$ we
have
\[
h\boxtimes(1,1)=\KS(h).
\]
\end{satz}

\begin{proof}
Verbatim the mirror image of Theorem~\ref{satz:wertseite}, with
Lemma~\ref{lem:sigmamono} in place of the $\mu$-monotonicity. The
homomorphism property of $\KS$: for $\boxplus$,
$\max(\sigma P,\sigma Q,\sigma(P{+}Q))=\sigma(P{+}Q)$ by the
lemma; for $\boxtimes$ both sides are
$(\sigma(PQ),PQ)$.
\end{proof}

\begin{bemerkung}
It is \emph{the same} Hori number $(1,1)$ --- the one in its
smallest home --- that on the value side carries $\KV$ and on the
structure side $\KS$ as inner multiplication. The non-commuting of
the compactifications is the non-commuting of two multiplications
by one and the same element. Incidentally, the transformation
monoid generated by $\KS,\KV$ is finite: it has exactly six
elements with the relation
$\KV\KS\KV=\KS\KV$ (via a stability lemma: $\KS$ preserves
value-compactness) and five idempotents
(proof in Appendix~\ref{app:monoid}).
\end{bemerkung}

\section{Third main result: the side-fidelity theorem}

Call an operation on $\Hori$ \emph{value-carried} if the result
value is the sum resp.\ product of the operand values (the result
horizon is free, provided it is valid), and
\emph{structure-carried} if the result polynomial is the sum resp.\
product of the operand polynomials (horizon free $\ge\sigma$).
Theorems~\ref{satz:wertseite} and~\ref{satz:dualitaet} show: both
\emph{side-pure} combinations are law-abiding with suitable horizon
rules. The mixed ones never are:

\begin{satz}[side fidelity]\label{satz:seitentreue}
\emph{(A)} If $\oplus$ is value-carried and $\otimes$
structure-carried, then already left-distributivity is
unsatisfiable for every choice of horizon rules.
\emph{(B)} If $\oplus$ is structure-carried and $\otimes$
value-carried, then left-distributivity and commutativity of
$\otimes$ are jointly unsatisfiable.
\end{satz}

\begin{proof}
\emph{(A).} Let $u=(1,1)$ (structure $1$) and $e_1=u$,
$e_m=e_{m-1}\oplus u$, so $\val(e_m)=m$; the structure $Q_m$ of
$e_m$ is some valid representation of $m$.
Left-distributivity yields inductively
\begin{equation}\label{eq:kette}
\val(a\otimes e_m)=m\cdot\val(a\otimes u).
\end{equation}

\emph{Step 1: every $Q_m$ is constant.} For $c\ge1$, $a_c=(c,c)$
has the constant structure $c$; constants evaluate to themselves
everywhere, so $\val(a_c\otimes u)=c$. With \eqref{eq:kette} it
follows that $c\,Q_m(y_c{+}1)=mc$, hence $Q_m(y_c{+}1)=m$, where
$y_c\ge\sigma(cQ_m)=\max_d(d+c\,q_d)\ge c$ grows without bound. All
evaluation points lie in the strictly monotone range of $Q_m$
(arguments $\ge\sigma+1\ge\deg+2$); for $\deg Q_m\ge1$ there would
be at most one such point. Hence $Q_m$ is constant, and with
$a=u$ in \eqref{eq:kette} it follows that $Q_m=m$.

\emph{Step 2: constants die.} Let $g=(2,3)$ with structure
$X$ and $r=\val(g\otimes u)=t_g{+}1\ge3$ fixed. By Step~1,
$g\otimes e_m$ has the structure $mX$ with $\sigma(mX)=m{+}1$, so
$\val(g\otimes e_m)=m(T{+}1)\ge m(m{+}2)$; by \eqref{eq:kette}
this value is $mr$ --- a contradiction for $m>r-2$.

\emph{(B).} Now structures add and values multiply. Let $f_1=u$,
$f_m=f_{m-1}\oplus u$; then $f_m$ has the constant structure $m$
and the value $m$. Let $A$ be the structure of
$g\otimes u$ (a representation of $3$). Inductively, $g\otimes f_m$
has the structure $mA$; as an $\otimes$-result of value $3m$,
realized on its horizon $T_m$, validity (the digit bound) forces
$T_m\ge\sigma(mA)\ge m$ and $A(T_m{+}1)=3$. As in Step~1,
$A$ is constant, so $A=3$. Now consider $f_2'=g\oplus g$ with
structure $2X$, horizon $H\ge\sigma(2X)=3$, and value
$2(H{+}1)\ge8$. Left-distributivity and commutativity give
\[
u\otimes(g\oplus g)=(g\otimes u)\oplus(g\otimes u),
\]
on the right with structure $2\cdot3=6$, hence value $6$; on the
left stands the value $2(H{+}1)\ge8$. Contradiction.
\end{proof}

\begin{bemerkung}
Both contradictions use the same property:
$\sigma(m\cdot P)$ grows linearly in $m$, because the digit bound
$c_d\le n-d$ forces large coefficients into large horizons. The
rigidity of the two identities is thus a direct consequence of the
defining digit condition $0\le a_i\le i$. The non-commutative mixed
case of type~(B) is likewise excluded under two-sided
distributivity: right-distributivity mirrors the chain
construction and replaces commutativity. And side fidelity is by
now \emph{three-sided}: the fraction identity (Cantor coordinate)
carries no arithmetic of its own -- its addition can exceed one and
is never total -- and cannot be interwoven with either of the two
sides (value-$\oplus$ fails at the constant factorial quotient,
structure-$\oplus$ at the half-coefficient obstruction
$2A(Y)=Y^{\underline k}$; proofs in
Appendix~\ref{app:side}).
There are exactly two computational worlds.
\end{bemerkung}

\section{The depth law of the commutator}

The defect
$\delta(h)=\operatorname{hor}(\KS\KV h)-\operatorname{hor}(\KV\KS h)$
is unbounded on both sides, but extremely asymmetric: upwards, $+2$
first occurs at $h=(9,720)$ ($720=6!=10\cdot9\cdot8$ has in
$H_9$ the structure $\ff X3$; secured exhaustively); the further
thresholds are bounded only from above: $+3$ at the latest from
$n=14$, $+4$ at the latest from $n=20$, $+5$ at the latest from
$n=26$. Downwards, the depth $-g$ is tied to the \emph{shell gap};
the jump positions are
\[
m_g=3,\ 15,\ 84,\ 533,\ 3883,\ 31998,\ 294875\qquad(g=1,\dots,7).
\]

\begin{satz}\label{satz:asymptotik}
With $c_g=(g{+}2)!\sum_{k\ge1}k/(g{+}k{+}1)!=1+O(1/g)$ we have
\[
\frac{(g+2)!}{c_g}\le m_g+1\le(g+2)!,
\qquad\text{in particular}\quad
\frac{m_g}{(g+2)!}\to1 .
\]
\end{satz}

\begin{proof}
Reindexing yields exactly
$M_{m-g}(m{+}1)/m!=(m{+}1)\sum_{k=1}^{m-g}k/(g{+}k{+}1)!$, where
$M_s=\sum_{d<s}(s{-}d)\ff Xd$ is the maximal polynomial. The
right-hand side grows strictly in $m$; for $m{+}1\ge(g{+}2)!$ the
term $k{=}1$ already suffices, and for $m{+}1<(g{+}2)!/c_g$ even
the full series is too small. The bound
$c_g\le1+\sum_{k\ge2}k(g{+}3)^{1-k}$
gives $c_g=1+O(1/g)$.
\end{proof}

\begin{satz}[depth law, exact equality]\label{satz:tiefengesetz}
For every shell $m\ge3$,
$\min\{\delta(h)\mid\mu(\val(h))=m\}=-g(m)$ with
$g(m)=m-s_{\min}(m)$,
$s_{\min}(m)=\min\{s\mid M_s(m{+}1)\ge m!\}$; the terminal
representation of $x^*=M_{s_{\min}}(m{+}1)$ is a witness.
\end{satz}

\begin{proof}
\emph{Bound.} Let $h=(n,x)$ with $\mu(x)=m$, so $x\ge m!$.
Left: $\KV(h)=(m,x)$, and for $P=\operatorname{struct}(x,m)$ we
have $P\preceq M_{\sigma(P)}$ coefficientwise, hence
$x=P(m{+}1)\le M_{\sigma(P)}(m{+}1)$; since $x\ge m!$ it follows
that $\sigma(P)\ge s_{\min}$ and thus
$\hor(\KS\KV(h))=\sigma(P)\ge s_{\min}$. Right: for
$P_n=\operatorname{struct}(x,n)$ the evaluation is monotone
(nonnegative coefficients), hence
$\val(\KS(h))=P_n(\sigma(P_n){+}1)\le P_n(n{+}1)=x<(m{+}1)!$ and
thus $\hor(\KV\KS(h))\le m$. Together
$\delta(h)\ge s_{\min}-m=-g(m)$.

\emph{Attainment.} We have
$x^*=M_{s_{\min}}(m{+}1)\in[m!,(m{+}1)!)$
-- the lower bound by the definition of $s_{\min}$, the upper
since $M_{s_{\min}}\preceq M_m$ and $M_m(m{+}1)=(m{+}1)!-1$. Let
$h^*=(x^*,x^*)$ be the terminal representation. $\KS$ fixes $h^*$
(constant shape), hence
$\hor(\KV\KS(h^*))=\mu(x^*)=m$. On the other hand, $M_{s_{\min}}$
is a valid digit sequence in horizon $m$
($c_d\le s_{\min}-d\le m-d$) with value $x^*$, hence by the
bijection from Theorem~\ref{satz:bij} \emph{the} digit
sequence: $\operatorname{struct}(x^*,m)=M_{s_{\min}}$, and
$\hor(\KS\KV(h^*))=\sigma(M_{s_{\min}})=s_{\min}$. Thus
$\delta(h^*)=s_{\min}-m=-g(m)$.
\end{proof}

Empirically, $(g{+}2)!/c_g$ even equals, to within $\pm1$, the
value $m_g$; the prediction $m_7=294\,875$ derived from this was
confirmed \emph{exactly}, as were two earlier predictions of this
law ($-3$ from shell ${\approx}100$; $m_6=31998$ against a
predicted ${\approx}31930$).

\section{The horizon calculus}

For a $k$-ary operator $T$ on the structure semiring
$\Struk$ let
\[
\Gamma_T(r_1,\ldots,r_k)=\max\{\sigma(T(P_1,\ldots,P_k)):
\sigma(P_i)\le r_i\}
\]
be the maximal structure horizon of the result.
We write $P\preceq Q$ for the coefficientwise order in the falling
factorial basis and
$M_r(X)=\sum_{d<r}(r{-}d)\,\ff Xd$ for the maximally filled
polynomial with $\sigma(M_r)=r$.

\begin{satz}[extremal principle]\label{satz:extremal}
If $T$ is coefficient-monotone in each argument --- which holds for
every operator composed of addition, multiplication, and $\Delta$,
since all structure constants are non-negative ---, then
\[
\Gamma_T(r_1,\ldots,r_k)=\sigma\bigl(T(M_{r_1},\ldots,M_{r_k})\bigr).
\]
\end{satz}

\begin{proof}
$\sigma(P)\le r$ is equivalent to $P\preceq M_r$ (digit bound
$c_d\le r-d$), and $\sigma$ is $\preceq$-monotone; monotonicity of
$T$ yields the upper bound, $M_{r_i}\in\mathcal F_{r_i}$ the
attainment.
\end{proof}

Thus every basic operation possesses its own exact growth law:
$\Gamma_+(r,s)=r+s$ (linear, sharp at constants),
$\Gamma_\times(r,s)=\beta(r,s)=\sigma(M_rM_s)$ (essentially
factorial, see below) --- and for the finite difference:

\begin{satz}[difference overflow]\label{satz:differenzueberlauf}
For $r\ge2$ we have
$\Gamma_\Delta(r)=\lfloor(r+1)^2/4\rfloor-1$, attained at the
single-digit structure $P=(r{-}d)\,\ff Xd$, $d=\lfloor(r+1)/2\rfloor$.
\end{satz}

\begin{proof}
By Theorem~\ref{satz:extremal},
$\Gamma_\Delta(r)=\sigma(\Delta M_r)$ with coefficients
$b_k=(k{+}1)(r{-}k{-}1)$, hence
$\Gamma_\Delta(r)=\max_{0\le k\le r-2}(k{+}1)(r{-}k)-1$; the
product of two numbers with sum $r{+}1$ is at most
$\lfloor(r{+}1)^2/4\rfloor$. (Confirmed exhaustively up to $r=8$:
$1,3,5,8,11,15,19$.)
\end{proof}

So $\Delta$ lowers the degree and raises the capacity
quadratically: the degree measures dimension, $\sigma$ measures
capacity; a degree filtration does not see this effect.

The calculus also captures substitution:

\begin{satz}[composition closure]\label{satz:komposition}
$\Struk$ is closed under composition, and $\circ$ is
coefficient-monotone in both arguments; in particular
$\Gamma_\circ(r,s)=\sigma(M_r\circ M_s)$.
\end{satz}

\begin{proof}[Proof sketch]
Because $P\circ Q=\sum_d c_d(P)(Q)^{\underline d}$, it suffices
that $(Q)^{\underline d}\in\Struk$. Interpreting $Q(x)$ as the
number of pairs consisting of a template of arity $k$ ($a_k$ of
each kind) and an injection $[k]\hookrightarrow[x]$, one sees that
$Q(x)^{\underline d}$ counts the $d$-tuples of pairwise distinct
such pairs. Classification by the union of images yields
$Q(x)^{\underline d}=\sum_m N_m\binom xm$; on the exhausting
tuples, $S_m$ acts \emph{freely} by relabelling the ground set (a
permutation that fixes every injection fixes the exhausted union
pointwise), so $m!\mid N_m$ and
$(Q)^{\underline d}=\sum_m(N_m/m!)X^{\underline m}\in\Struk$ ---
positivity and integrality in one stroke. Monotonicity in $Q$
follows because larger template supplies only enlarge the tuple
sets.
(Confirmed exhaustively for $r,s\le4$; the complete proof is in
Appendix~\ref{app:calc}.)
\end{proof}

The diagonal $\Gamma_\circ(r,r)=1,4,23,6541,477149751,\ldots$
grows tower-like and is the ninth OEIS-unknown sequence of the
project; its asymptotics are open. The shift operator
$E\,P(X)=P(X{+}1)$, as the simplest special case, costs exactly
$\Gamma_E(r)=r+\lfloor r^2/4\rfloor=\Gamma_\Delta(r{+}1)$ ---
shifting costs as much as differencing one horizon higher
(Appendix~\ref{app:calc}); this
value sequence is the classical quarter-squares sequence (A024206),
new is the law, not the sequence. The continuous analogue of the
composition closure is classical (absolutely monotone functions,
Faà di Bruno); Theorem~\ref{satz:komposition} is the discrete
lattice sharpening with the additional divisibility $m!\mid N_m$.
Additively,
$(\mathcal F_r)_r$ with $\mathcal F_r=\{P:\sigma(P)\le r\}$,
$|\mathcal F_r|=(r{+}1)!$, is a classically compatible filtration;
multiplicatively it is non-linear with optimal growth law
$\beta$. Via the Newton expansion ($c_d=\Delta^dP(0)/d!$),
$\sigma(P)=\max_{\Delta^dP(0)>0}\bigl(d+\Delta^dP(0)/d!\bigr)$ is a
basis-free \emph{capacity height}; it is not a bit complexity
($\sigma$ of a constant $c$ is $c$ itself).

\section{The counting width: a \#P interpretation}

Because $X^{\underline d}=d!\binom Xd$, every $P\in\Struk$ has
non-negative integer coefficients in the binomial basis
($b_d=c_d\,d!$) and thus lies in the class of
\emph{binomial-good} polynomials, which are characterized as the
(relativizing) polynomial \#P closure operations
(Ikenmeyer--Pak, FOCS 2022, arXiv:2204.13149). $\Struk$ is a
\emph{proper} subclass ($\binom X2\notin\Struk$); the
$d!$-divisibility corresponds to ordered tuples instead of
unordered selections.

\begin{satz}[Hori \#P theorem]\label{satz:sharpp}
For $f\in\#\mathrm P$ and $P\in\Struk$ we have
$P\circ f\in\#\mathrm P$.
\end{satz}

\begin{proof}
Guess $d\le\deg P$, a variant $\ell\le c_d$, and $d$ witness
candidates; accept upon $d$ pairwise distinct valid witnesses. The
number of accepting paths is
$\sum_d c_d f(x)^{\underline d}=P(f(x))$; since $P$ is fixed,
everything remains polynomial.
\end{proof}

In this reading the Hori operations acquire a counting semantics:
addition is disjoint union of counting constructions,
multiplication couples constructions (the linearization
coefficients $\binom pj\binom qj j!$ count the overlap patterns),
$\Delta$ counts the structures enabled by one additional witness
($(n{+}1)^{\underline d}-n^{\underline d}=d\,n^{\underline{d-1}}$),
and $\sigma(P)=\max(d+c_d)$ measures the maximal combination of
witness count and variant count of a building block
(\emph{counting width}; geometrically the smallest staircase above
the coefficient profile). Each $\mathcal F_r$ is thus a finite
family of exactly $(r{+}1)!$ \#P closure operators, and the horizon
calculus quantifies optimally the growth of the counting width
under $+$, $\times$, $\Delta$, and $\circ$. Demarcation: the
binomial basis conjecture of Ikenmeyer--Pak (whose univariate
version would imply P$\,\ne\,$NP) remains untouched; the Hori
operators lie on the known, relativizing side.
New is exclusively the quantitative hierarchy.

\section{Nine counting sequences}

The following sequences originate from Horimetrik and produced no
hit in the OEIS (last checked on 14~July 2026; a negative finding
is no proof of novelty). Submission is under way; extended data
tables (jump positions up to $g=20$, jumper, non-jumper, and
fixed-point sequences up to $n=40$, interchange up to
$H_{12}$, 78 single-digit terms) are available as b-files in the
repository:

\begin{center}
\begin{tabular}{ll}
\toprule
Sequence & First terms\\
\midrule
$\beta(r,r)$, maximal structure height of products &
$1,6,22,87,407,2473,19527,\dots$\\
Defect $D(n)$ (one-step shell jumpers) &
$1,8,60,360,2520,21168,211680$\\
two-sided fixed points $F(n)=n\cdot n!-D(n)$ &
$4,17,88,540,3960,32760$\\
jump positions $m_g$ of the depth law &
$3,15,84,533,3883,31998,294875$\\
interchange maxima $\max\Omega$ on $H_n$ &
$0,1,6,43,351,2873,26443,270291$\\
interchange zero count $\#\{\Omega=0\}$ &
$2,5,13,23,61,111,274,564$\\
non-jumpers $n!-D(n)$ &
$5,16,60,360,2520,19152$\\
single-digit factorials $\{n: n!\ \text{single-digit in } H_{\mu(n!)}\}$ &
$1,2,5,7,11,23,29,39,59,119$\\
composition height $\Gamma_\circ(r,r)$ &
$1,4,23,6541,477149751,\dots$\\
\bottomrule
\end{tabular}
\end{center}

Here $\Omega(h)=\val(LZh)-\val(ZLh)$ is the defect between the two
growth orders; the positivity $\Omega\ge0$ is proved
(multiplication lemma with a sharp regime boundary; additionally
co-confirmed exhaustively up to $H_{12}$, i.e.\ over
$6.2\cdot10^9$ elements), the
equality cases are characterized and completely unfolded into digit
conditions (two inequalities per level with growing excess), and
the zero count obeys a counting formula derived from this (exact
for all measured terms) with a proved upper bound
$(n{+}1)^{n+1}/n!$ — indifference to the growth order is
asymptotically extremely rare
(proofs in Appendix~\ref{app:interchange}). For
$\beta(r,r)$ the asymptotics are proved:
$\ln(\beta(r,r)/r!)=2\sqrt r+O(\log r)$ (reduction to one maximal
term plus elementary bounds; Appendix~\ref{app:beta}). By the shell-jump lemma, $D(n)$
counts exactly the elements that a single leading zero lifts into
the next factorial shell; the corresponding digit formula runs over
the eigenrepresentation of $n!$, whose single-digit cases are
completely characterized ($n=(i{+}1)!/d-1$, $1\le d\le i$). All
nine sequences will be submitted to the OEIS; the A-numbers will be
added to this text once assigned, so that text and database
reference each other.

\section{Context and honesty}

The building blocks are known mathematics: mixed-radix systems and
factoradics, Lehmer codes, falling factorials and finite
differences, direct limits, product semirings. The closest known
neighbours of the zero semantics are abstract numeration systems;
there, leading zeros change the value via the radix order, without
a semiring theory being erected on top. What we have not found in
this form is the constellation: finite factorial horizons with
running evaluation, value/structure duality with two projectors to
the same one, and a rigidity theorem that forbids mixing the
identities.

The philosophical motivation --- infinity as a staircase of finite
horizons --- is motivation, not argument. The robust version reads:
infinite structures can be understood as systems of finite
horizons, and \emph{different notions of consistency between the
horizons generate different global spaces}
(Theorem~\ref{satz:limits}).

The base space itself captures the situation most precisely. If one
augments the digits by a leading zero, $e=(0,a_1,\dots,a_n)$, then
$0\le e_j<j$: as a set, $H_n$ is the space of inversion sequences
of length $n{+}1$, that is, $H_n\cong S_{n+1}$. The structure
horizon becomes the permutation statistic
$\sigma(P_a)=n-\min_{e_j>0}((j{-}1)-e_j)$ with exact distribution
$|\{\sigma=r\}|=r\cdot r!$ (Appendix~\ref{app:inv}); whether it corresponds to a classical
statistic is open. Coordinatewise, each horizon becomes a
distributive lattice (the \emph{middle order} on permutations,
arXiv:2405.08943), in which $\mathcal F_r$ is the principal ideal
below $M_r$. The inhabitants are thus known permutation codes; new
are the operations on them --- the reversed valuation, the horizon
transitions, the duality, the side fidelity, and the capacity
calculus. That is the most favourable situation for a young theory:
known base space, new theorems.

Several earlier conjectures of this project were refuted by our own
computations (one bound conjecture fell twice to test ranges that
were too small; one formula conjecture at its fourth data point).
We regard this as a quality feature of the method, not as a
weakness of the subject.

\section{Open questions}

(1)~The \emph{upper} thresholds of the depth law exactly (so far
only ``at the latest from'' bounds; the lower side is completely
proved). (2)~Fine structure of the $\beta$-asymptotics (conjectured:
$-\tfrac32\ln r$ in the logarithm term). (3)~General digit
statistics of $n!$ in its own shell (the single-digit cases are
completely characterized); from this, the distribution asymptotics
of $D(n)$. (4)~Classification of the \emph{coupled} arithmetics
within one side (beyond the product constructions).
(5)~A lower bound for the interchange zero count.
(6)~$\Gamma_T$ for further operators (composition, shift) and the
question whether the capacity height $\max_d(d+c_d)$ has already
been studied elsewhere. (7)~Formalization of the side-fidelity
theorem (the falling factorials are already available in
Lean~4/Mathlib).

\bigskip
\noindent\textbf{Reproducibility.} All enumerations and
verifications of this text are available as a Python reference
implementation with self-tests; every result carries in the working
files the status proved/conjectured/refuted with a date. Source
code, b-files, and the detailed book (German and English, with all
proofs) are available at
\texttt{github.com/}\allowbreak\texttt{christianbuerckert/horimetrik}.


\section*{Acknowledgements}
The mathematics of this text was developed in close collaboration
with Claude (Anthropic); GPT (OpenAI) contributed ideas and early
literature scouting. All results were verified computationally
against a reference implementation with self-tests, and the
manuscript passed independent adversarial review passes before
release. Responsibility for the content rests with the author.

\appendix

\section{Proofs for the horizon calculus}\label{app:calc}

\emph{Complete proof of Theorem~\ref{satz:komposition}
(composition closure).}

\begin{proof}
Since $P\circ Q=\sum_d c_d(P)\,(Q)^{\underline d}$ with
$(Q)^{\underline d}=Q(Q{-}1)\cdots(Q{-}d{+}1)$, it suffices to
show that $(Q)^{\underline d}$ is an admissible shape. Write
$Q=\sum_k a_k\ff Xk$ and choose pairwise disjoint sets $T_k$ with
$|T_k|=a_k$ (\emph{templates of arity $k$}). Let a
\emph{$Q$-structure} on $[x]$ be a pair $(t,\varphi)$ consisting
of a template $t\in T_k$ and an injection
$\varphi\colon[k]\hookrightarrow[x]$; their number is
$\sum_k a_k\,\ff xk=Q(x)$. Hence $Q(x)^{\underline d}$ counts the
$d$-tuples of pairwise distinct $Q$-structures on $[x]$.

Classify the tuples by the union $U\subseteq[x]$ of the images of
all injections involved. The number $N_m$ of tuples whose image
union entirely exhausts a fixed $m$-set depends only on $m$, and
for all $x\in\NN$ the polynomial identity
$Q(x)^{\underline d}=\sum_m N_m\binom xm$ holds. The symmetric
group $S_m$ acts on the exhausting tuples by renaming the ground
set, $\pi\cdot(t,\varphi)=(t,\pi\circ\varphi)$, and this action is
\emph{free}: if $\pi$ fixes every member of the tuple, then $\pi$
fixes every image pointwise, hence by exhaustion all of $[m]$,
i.e.\ $\pi=\mathrm{id}$. (Templates of arity 0 do no harm: they
contribute nothing to the union.) Thus $m!\mid N_m$, and
$(Q)^{\underline d}=\sum_m(N_m/m!)\,\ff Xm$ has coefficients in
$\NN$ -- positivity and integrality in one stroke.

Monotonicity in $P$: an $\NN$-linear combination of the
$(Q)^{\underline d}$. Monotonicity in $Q$: $a_k\le a_k'$ yields
$T_k\subseteq T_k'$; every tuple of $Q$-structures is one of
$Q'$-structures with the same image union, hence $N_m\le N_m'$
term by term. (Closure and extremal principle confirmed
computationally and exhaustively for $r,s\le4$, the closure
additionally for 300 random pairs up to $\sigma=8$.)
\end{proof}

\begin{satz}[shift overflow]
For $E\,P(X)=P(X{+}1)$,
$\Gamma_E(r)=r+\lfloor r^2/4\rfloor=\Gamma_\Delta(r{+}1)$.
\end{satz}

\begin{proof}
$E$ is linear with nonnegative constants
($\ff{(X{+}1)}d=\ff Xd+d\,\ff X{d-1}$), so the extremal principle
applies: $\Gamma_E(r)=\sigma(E\,M_r)$ with coefficients
$c_k=(r{-}k)+(k{+}1)(r{-}k{-}1)$ for $k\le r{-}2$ and $c_{r-1}=1$;
thus $k+c_k=r+(k{+}1)(r{-}k{-}1)$, and the product of two numbers
with sum $r$ is at most $\lfloor r^2/4\rfloor$. The second
equality is the identity
$r+\lfloor r^2/4\rfloor=\lfloor(r{+}2)^2/4\rfloor-1$.
\end{proof}


\section{Product principle and the compactification
monoid}\label{app:monoid}

\begin{satz}[product principle]
Every commutative semiring structure $(\uplus,\odot)$ on the
excesses yields, via
$(x,\varepsilon)\oplus^{*}(y,\delta)=(x{+}y,\varepsilon\uplus\delta)$,
$(x,\varepsilon)\otimes^{*}(y,\delta)=(xy,\varepsilon\odot\delta)$,
a law-abiding arithmetic on $\Hori$.
\end{satz}

\begin{proof}
The map $(n,x)\mapsto(x,\,n-\mu(x))$ is a bijection
$\Hori\to\NN\times\NN$ (the condition $n\ge\mu(x)$ is exactly
$\varepsilon\ge0$; the inverse is
$(x,\varepsilon)\mapsto(\mu(x){+}\varepsilon,x)$). The operations
are defined componentwise, with $(\NN,+,\cdot)$ in the value
factor; a product of two commutative semirings satisfies all laws
componentwise. The selectors are valid, since the resulting
horizon is $\mu(\text{value})+\text{excess}\ge\mu$.
\end{proof}


\begin{satz}[compactification monoid]
$\KS$ preserves value-compactness. Consequently
$\KV\KS\KV=\KS\KV$, and the transformation monoid generated by
$\KS,\KV$ has exactly six elements -- $\mathrm{id}$,
$\KS$, $\KV$, $\KS\KV$, $\KV\KS$ and $\KS\KV\KS$ -- of which
all except $\KV\KS$ are idempotent.
\end{satz}

\begin{proof}
\emph{Stability:} let $(m,P)$ be value-compact; the element
$(0,0)$ is fixed by $\KS$, so let $m\ge1$ and hence
$P(m{+}1)\ge m!$. Let $s=\sigma(P)$; suppose we had
$P(s{+}1)<s!$. For $s\le t<m$, each term $c_d\ff Xd$ grows from
$t{+}1$ to $t{+}2$ by the factor $(t{+}2)/(t{+}2{-}d)$, hence,
since $d\le s{-}1\le t{-}1$, by at most $(t{+}2)/3$; thus
$P(t{+}2)<t!\,(t{+}2)/3\le(t{+}1)!$, and induction yields
$P(m{+}1)<m!$ -- contradiction.
\emph{Relation:} for every $h$, $\KV(h)$ is value-compact, by
stability so is $\KS(\KV(h))$, hence $\KV$ fixes this element.
\emph{Six elements:} every alternating word of length $\ge4$
contains $\KV\KS\KV$ and reduces; squares reduce by idempotence
-- at most six. Pairwise distinctness by explicit separating
witnesses (finite enumeration); in particular $h=(6,6)$
separates: $\KV\KS(h)=(3,6)$, but $\KS\KV\KS(h)=(2,5)$. Hence
$(\KV\KS)^2=\KS\KV\KS\ne\KV\KS$: not idempotent; the idempotence
of the remaining five is verified using the relation.
\end{proof}


\section{Side fidelity: the noncommutative case and the fraction
identity}\label{app:side}

\begin{satz}[noncommutative case]
If $\oplus$ is structure-carried and $\otimes$ value-carried and
both distributive laws hold, then the system is unsatisfiable --
without commutativity, associativity, or neutral elements.
\end{satz}

\begin{proof}
Right distributivity yields inductively
$P_{f_m\otimes a}=m\cdot P_{u\otimes a}$; since $f_m\otimes a$ has
the value $m\cdot\val(a)$ and the horizon
$\ge\sigma(m\,P_{u\otimes a})\ge m$ grows, $P_{u\otimes a}$ is
constant equal to $\val(a)$ (monotonicity argument as above).
In particular $P_{u\otimes g}=3$; the $g\oplus g$ finale then runs
word for word as in the commutative case via left distributivity.
\end{proof}


\begin{satz}[trilogy]
Fraction-carried addition is never total; value-carried addition
with fraction-carried multiplication, as well as
structure-carried addition with fraction-carried multiplication,
are unsatisfiable for every choice of selectors (left
distributivity suffices).
\end{satz}

\begin{proof}
\emph{Totality:} $F(u)+F(u)=\tfrac12+\tfrac12=1\notin[0,1)$, and
every element has fraction $<1$.

\emph{Value $\oplus$, fraction $\otimes$:} chains $e_m$ with
$\val(e_m)=m$, horizon $N_m\ge\mu(m)$, fraction $m/(N_m{+}1)!$.
Distributivity gives $\val(u\otimes e_m)=m\,w$ with
$w=\val(u\otimes u)$ -- an integer and $\ge1$, since the fraction
of $u\otimes u$ is $\tfrac14>0$. On the other hand, with $h_m$ the
horizon of $u\otimes e_m$,
$\val(u\otimes e_m)=\tfrac12\cdot\frac{m}{(N_m+1)!}(h_m{+}1)!$,
hence $(h_m{+}1)!/(N_m{+}1)!=2w>1$ constant -- but a
factorial quotient $>1$ is at least $N_m{+}2\to\infty$.

\emph{Structure $\oplus$, fraction $\otimes$:} chains $f_m$ with
constant structure $m$, horizon $N_m\ge m$, fraction
$m/(N_m{+}1)!$. Distributivity gives $P_{u\otimes f_m}=m\,A$ with
$A=P_{u\otimes u}$, $\deg A=:d$; comparing fractions on the
horizon $T_m$ of $u\otimes f_m$ yields, after cancelling $m$,
\[
2\,A(T_m{+}1)=(T_m{+}1)^{\underline{\,k_m}},\qquad
k_m=T_m-N_m\ge1 .
\]
The product on the right is at least $(N_m+k_m/2)^{k_m/2}$, the
left-hand side at most $C(N_m{+}k_m{+}1)^d$; since $N_m\to\infty$,
$k_m\le2d{+}1$ for large $m$, some value $k^*$ occurs infinitely
often, and $2A(Y)=Y^{\underline{k^*}}$ holds as a polynomial
identity. Then $A$ would have the coefficient $\tfrac12$ --
digits are integers. (General multipliers $a$: the same
computation with $c=(n_a{+}1)!/v_a>1$ instead of $2$.)
\end{proof}


\section{The interchange defect: positivity, sharpness, zero
count}\label{app:interchange}


For $x<(n{+}1)!$ let $z_n(x):=P(n{+}2)$ with
$P=\operatorname{struct}(x,n)$ -- the value after one zero
extension. The \emph{interchange defect} of an element
$h=(n,x)$ is
$\Omega(h)=z_n(x)-z_{n+1}(x)$: the value after ``grow first, then
lift value-preservingly'' minus the value after ``lift first, then
grow''.

\begin{lemma}[recursion formula]
For $x=q(n{+}1)+r$, $0\le r\le n$:
$z_n(x)=r+(n{+}2)\,z_{n-1}(q)$.
\end{lemma}

\begin{proof}
From the digit recursion it follows that
$P_{n,x}(X)=r+X\cdot P_{n-1,q}(X{-}1)$, since the coefficient
shift by one degree is, by
$\ff Xd=X\cdot\ff{(X{-}1)}{d-1}$, multiplication by $X$ with a
shift of argument. Evaluate at $X=n{+}2$.
\end{proof}

\begin{lemma}[strict monotonicity and superadditivity]
$z_n(x)<z_n(x{+}1)$, consequently $z_n(y{+}s)\ge z_n(y)+s$.
\end{lemma}

\begin{proof}
Induction on $n$; base $n=1$: $z_1(0)=0<1=z_1(1)$. Without a
carry, $z_n$ grows by exactly 1; with a
carry ($r=n$),
$z_n(x{+}1)-z_n(x)=(n{+}2)(z_{n-1}(q{+}1)-z_{n-1}(q))-n
\ge(n{+}2)-n=2$.
\end{proof}

\begin{lemma}[multiplication lemma]
For $c\ge\ell{+}2$ and $cy<(\ell{+}1)!$,
$z_\ell(cy)\ge(c{+}1)\,z_\ell(y)$.
\end{lemma}

\begin{proof}
Induction on $\ell$; the base cases are empty and trivial,
respectively. Let
$y=q(\ell{+}1)+r$ and $cr=s(\ell{+}1)+u$ with $0\le u\le\ell$;
then $cy=(cq{+}s)(\ell{+}1)+u$ and, by the recursion formula,
$z_\ell(cy)=u+(\ell{+}2)\,z_{\ell-1}(cq{+}s)$. Superadditivity
and the induction hypothesis (the regime is preserved:
$c\ge\ell{+}2>(\ell{-}1){+}2$; validity: $cq<\ell!$) give
$z_{\ell-1}(cq{+}s)\ge(c{+}1)z_{\ell-1}(q)+s$. It remains to show
$u+(\ell{+}2)s\ge(c{+}1)r$: with $c=\ell{+}1{+}e$, $e\ge1$, we
have $u=er\bmod(\ell{+}1)$, and the inequality is equivalent to
$r(c{-}\ell{-}1)=er\ge er\bmod(\ell{+}1)=u$ -- true.
\end{proof}

\begin{satz}[interchange positivity]
$\Omega(h)\ge0$ for all $h\in\Hori$.
\end{satz}

\begin{proof}
By the definition of $\Omega$ it must be shown that
$z_{n+1}(x)\le z_n(x)$ for $x<(n{+}1)!$. Let
$x=q'(n{+}2)+r'$. Then
\[
z_{n+1}(x)=r'+(n{+}3)\,z_n(q')
\le r'+z_n\bigl((n{+}2)q'\bigr)
\le z_n\bigl(q'(n{+}2)+r'\bigr)=z_n(x),
\]
first by the multiplication lemma ($c=n{+}2$, exactly at the
regime boundary), then by superadditivity.
\end{proof}

\begin{satz}[equality cases]
$\Omega(h)=0$ holds if and only if both steps of this chain are
sharp: the multiplication lemma with $c=n{+}2$ for the
\emph{sharp quotient} $q'$, i.e.\
$z_n((n{+}2)q')=(n{+}3)\,z_n(q')$, and superadditivity
carry-free, i.e.\ $z_n((n{+}2)q'+r')=z_n((n{+}2)q')+r'$.
\end{satz}

\begin{proof}
$\Omega(h)$ is the sum of the two deviations
$z_n((n{+}2)q')-(n{+}3)z_n(q')$ and
$z_n((n{+}2)q'+r')-z_n((n{+}2)q')-r'$; both are nonnegative by
the multiplication lemma and superadditivity respectively, so
their sum is zero if and only if both are zero.
\end{proof}

\begin{satz}[digit unfolding of sharpness]
For $c\ge\ell+2$, $cy<(\ell+1)!$, $y=q(\ell+1)+r$, $e=c-\ell-1$
we have: $z_\ell(cy)=(c{+}1)z_\ell(y)$ if and only if $er\le\ell$,
$(cq\bmod\ell)+r\le\ell-1$, and the same condition holds
recursively for
$(\ell{-}1,c,q)$ (termination: for $\ell=0$ the condition is
empty).
\end{satz}

\begin{proof}
The chain $z_\ell(cy)=u+(\ell+2)z_{\ell-1}(cq+s)$ from the proof
of the multiplication lemma is sharp throughout exactly then: the
arithmetic estimate when $er\le\ell$ (then $s=r$, $u=er$),
superadditivity when the $s=r$ single steps starting from $cq$
are carry-free -- their units digit is $cq\bmod\ell$, the
condition therefore $(cq\bmod\ell)+r\le\ell-1$ --, the third
recursively by definition. Verified exhaustively for $\ell\le7$.
\end{proof}

\begin{satz}[zero-count bound]
$\#\{\Omega=0\text{ on }H_n\}\le(n{+}1)^{n+1}/n!$, so the
proportion falls at least like $e^{-(1+o(1))n\ln n}$.
\end{satz}

\begin{proof}
By the equality theorem, every such element is determined by
$(q',r')$ with sharp $q'$, and carry-freeness forces
$r'\le n$ (the units digit grows by one per step and must stay
below its maximum $n$ before each step). By the digit unfolding,
sharpness forces at level $\ell$ the
digit bound $r_\ell\le\lfloor\ell/(n{+}1{-}\ell)\rfloor$, and
$\lfloor\ell/j\rfloor+1\le(n{+}1)/j$ with $j=n{+}1{-}\ell$ yields
the product $(n{+}1)^n/n!$ for the $q'$; together with the
$n{+}1$ possibilities for $r'$ the bound follows.
\end{proof}


\section{Digit formula, single-digit factorials, and the
$\beta$ asymptotics}\label{app:beta}

\begin{satz}[digit formula for the non-jumpers]
Let $d=(d_1,\ldots,d_n)$ be the $H_n$ digit sequence of $n!$ and
$P$ the first position with maximal digit ($d_P=P$), or $P=n$
if none exists. Then
\[
N(n)=\sum_{p=1}^{P}\min(d_p,p)\cdot\frac{n!}{p!}.
\]
\end{satz}

\begin{proof}
$N(n)$ counts digit sequences with $a_i\le i-1$ and value $<n!$.
We decompose by the first position $p$ at which $a$ differs from
$d$. The prefix must agree with $d$ -- possible only as long as
$d_j\le j-1$, i.e.\ up to and including the first maximal digit.
At position $p$ we need $a_p<d_p$ and $a_p\le p-1$:
$\min(d_p,p)$ possibilities. The suffix is free within the caps:
$\prod_{j>p}j=n!/p!$ possibilities, and every such suffix stays
strictly below, because the full (uncapped) suffix space fills
the interval $[0,w_p)$ with the place weight
$w_p=(n{+}1)!/(p{+}1)!$ bijectively. Every non-jumper is counted
exactly once.
\end{proof}


\begin{satz}[$\beta$ reduction]
$\ln\beta(r,r)=\ln\max_{t}\binom{r-1}{t}^2(r{-}1{-}t)!+O(\log r)$.
\end{satz}

\begin{proof}
Lower bound: the triple $(p,q,j)=(r{-}1,r{-}1,r{-}1{-}t)$
contributes to degree $r{-}1{+}t$ the summand
$\binom{r-1}{t}^2(r{-}1{-}t)!$, and all contributions are
nonnegative. Upper bound: each contribution is at most
$r^2\binom{r-1}{j}^2j!$, and per degree there are at most $r^3$
triples; hence $\beta\le2r+r^5\max_t(\cdot)$.
\end{proof}


\begin{satz}[$\beta$ asymptotics]
$\ln(\beta(r,r)/r!)=2\sqrt r+O(\log r)$.
\end{satz}

\begin{proof}
By the reduction theorem it suffices to prove the statement for
the maximal term
$\binom st^2(s-t)!/r!=s!/(r\,t!^2(s-t)!)$ with $s=r-1$. From
above: Stirling bounds give
$\ln(s!/(t!^2(s-t)!))\le-h(t)+O(\log s)$
with $h(t)=2t\ln t-t+(s{-}t)\ln(s{-}t)-s\ln s$; concavity yields
$h(t)\ge t(2\ln t-\ln s-2)$, and this function has its minimum
$-2\sqrt s$ at $t=\sqrt s$. From below: $t=\lceil\sqrt s\rceil$
with $\binom st\ge(s-t)^t/t!$ gives
$t\ln((s-t)/t^2)+2t-O(\log r)=2\sqrt s-O(\log r)$.
\end{proof}


\begin{satz}[single-digit factorials]
$n!$ is a single digit in $H_n$ if and only if $n=(i{+}1)!/d-1$
with $1\le d\le i$; the digit is $d$ at position $i$, and per
position there are exactly $i$ cases.
\end{satz}

\begin{proof}
A single digit $d$ at position $i$ means $n!=d\,(n+1)!/(i+1)!$,
hence $d=(i+1)!/(n+1)\le i$. Conversely, every $d\le i$ divides
$(i+1)!$, and $n+1=(i+1)!/d$ makes the equation exact; the
position exists, since $n{+}1=(i{+}1)!/d\ge(i{+}1)!/i\ge
i{+}1$ implies $i\le n$, and the unique digit expansion then
consists of exactly this digit.
$d$ runs bijectively through $1,\ldots,i$.
\end{proof}


\section{The inversion-sequence statistics}\label{app:inv}

\begin{satz}[inversion-sequence statistics]
$H_n$ is in bijection with the inversion sequences of length
$n{+}1$ and with $S_{n+1}$. The structure horizon is
$\sigma(P_a)=n-\min_{e_j>0}((j{-}1)-e_j)$ for $a\ne0$ (for $a=0$,
$\sigma=0$) with $e=(0,a_1,\ldots,a_n)$, and
$|\{a\in H_n:\sigma(P_a)=r\}|=r\cdot r!$ for $1\le r\le n$.
\end{satz}

\begin{proof}
Setting $e_1=0$ and $e_{i+1}=a_i$, we get $0\le e_j<j$ for all
$j=1,\ldots,n{+}1$ -- the defining condition of an inversion
sequence; the map $\pi\mapsto(e_j)_j$ with
$e_j=|\{i<j:\pi_i>\pi_j\}|$ is the classical bijection to
$S_{n+1}$, and both sides have $(n{+}1)!$ elements. For the
structure horizon: $c_d=a_{n-d}=e_{n-d+1}$, hence
$d+c_d=(n{+}1-j)+e_j$ with $j=n-d+1$; maximizing $d+c_d$ over
$c_d>0$ is minimizing $(j{-}1)-e_j$ over $e_j>0$, and
$\max(d+c_d)=n-\min_{e_j>0}((j{-}1)-e_j)$. Counting:
$\sigma(P_a)\le r$
means $a_i\le r-(n-i)$ for all $i$ with $a_i>0$, equivalently
$a\le M$ coordinatewise for the maximal profile belonging to
$H_r$; the number of such $a$ in $H_n$ is $(r{+}1)!$ (the digits
above position $r$ are forced to be zero, below they are free up
to the bound). Hence $|\{\sigma\le r\}|=(r{+}1)!$ and
$|\{\sigma=r\}|=(r{+}1)!-r!=r\cdot r!$. (Confirmed exhaustively
up to $n=6$, \texttt{experiments\_einordnung.py}.)
\end{proof}


\begin{thebibliography}{9}
\bibitem{gkp} R.~L. Graham, D.~E. Knuth, O.~Patashnik:
\emph{Concrete Mathematics}, 2nd ed., Addison--Wesley, 1994.
\bibitem{rigo} M.~Rigo, M.~Stipulanti: \emph{Revisiting regular
sequences in light of rational base numeration systems},
arXiv:2103.16966.
\bibitem{kiselman} G.~Kudryavtseva, V.~Mazorchuk: \emph{On
Kiselman's semigroup}, arXiv:math/0511374; O.~Ganyushkin,
V.~Mazorchuk: \emph{On Kiselman quotients of 0-Hecke monoids},
arXiv:1006.0316.
\bibitem{ybe} D.~Simon: \emph{Mixed radix numeration bases:
Horner's rule, Yang-Baxter equation and Furstenberg's conjecture},
arXiv:2405.19798.
\bibitem{ikenmeyerpak} C.~Ikenmeyer, I.~Pak: \emph{What is in \#P
and what is not?}, FOCS 2022, arXiv:2204.13149.
\bibitem{middleorder} M.~Bouvel, L.~Ferrari, B.~E.~Tenner:
\emph{Between weak and Bruhat: the middle order on permutations},
arXiv:2405.08943.
\end{thebibliography}

\end{document}
